\let\negmedspace\undefined
\let\negthickspace\undefined
\documentclass[journal]{IEEEtran}
\usepackage[a5paper, margin=10mm, onecolumn]{geometry}
%\usepackage{lmodern} % Ensure lmodern is loaded for pdflatex
\usepackage{tfrupee} % Include tfrupee package

\setlength{\headheight}{1cm} % Set the height of the header box
\setlength{\headsep}{0mm}     % Set the distance between the header box and the top of the text

\usepackage{gvv-book}
\usepackage{gvv}
\usepackage{cite}
\usepackage{amsmath,amssymb,amsfonts,amsthm}
\usepackage{algorithmic}
\usepackage{graphicx}
\usepackage{textcomp}
\usepackage{xcolor}
\usepackage{txfonts}
\usepackage{listings}
\usepackage{enumitem}
\usepackage{mathtools}
\usepackage{gensymb}
\usepackage{comment}
\usepackage[breaklinks=true]{hyperref}
\usepackage{tkz-euclide} 
\usepackage{listings}
% \usepackage{gvv}                                        
\def\inputGnumericTable{}                                 
\usepackage[latin1]{inputenc}                                
\usepackage{color}                                            
\usepackage{array}                                            
\usepackage{longtable}                                       
\usepackage{calc}                                             
\usepackage{multirow}                                         
\usepackage{hhline}                                           
\usepackage{ifthen}
\usepackage{lscape}
\begin{document}

\bibliographystyle{IEEEtran}



\title{2.10.85}
\author{EE25BTECH11057 - Rushil Shanmukha Srinivas
}
% \maketitle
% \newpage
% \bigskip
{\let\newpage\relax\maketitle}

\renewcommand{\thefigure}{\theenumi}
\renewcommand{\thetable}{\theenumi}
\setlength{\intextsep}{10pt} % Space between text and floats

\numberwithin{equation}{enumi}
\numberwithin{figure}{enumi}
\renewcommand{\thetable}{\theenumi}



\textbf{Problem:} Let P be the plane 3x+2y+3z=16 and let S: $\alpha\hat{i} + \beta\hat{j} + \gamma\hat{k} $ ,where $\alpha+\beta+\gamma=7$ and the distance of $(\alpha,\beta,\gamma)$ from the plane is $2/{\sqrt{22}}$ .Let $\vec{u},\vec{v},\vec{w}$ be three distinct vectors in S such that $|\vec{u}-\vec{v}|=|\vec{v}-\vec{w}|=|\vec{w}-\vec{u}|$. Let V be the volume of the parallelopiped determined by vectors $\vec{u},\vec{v},\vec{w}$. Then the value of (80/3)V is 

\textbf{Solution:}

\begin{align}
P:\; \vec{n}^\top \vec{x} = c,\qquad \vec{n}=\myvec {3\\2\\3}\;c=16.
\end{align}

The distance of point $\vec{P_0}=\myvec{\alpha\\ \beta\\ \gamma}$ from the plane is
\begin{align}
\operatorname{dist}(\vec{P_0},P)=\frac{|\vec{n}^\top \vec{P_0} - c|}{\norm{\vec{n}}}
\end{align}
Given $\alpha+\beta+\gamma=7$ and $\operatorname{dist}(\vec{P_0},P)=\tfrac{2}{\sqrt{22}}$, we have
\begin{align}
\frac{|3\alpha+2\beta+3\gamma-16|}{\sqrt{22}}=\frac{2}{\sqrt{22}}
\;\;\Longrightarrow\;\; \vec{n}^\top \vec{P_0}=18 \;\;\text{or}\;\; \vec{n}^\top \vec{P_0}=14.
\end{align}

Thus $S$ lies on the intersections
\begin{align}
\Pi:\; \vec{m}^\top \vec{x}=7,\qquad \vec{m}=\myvec{1\\1\\1},
\end{align}
with
\begin{align}
P_+: \vec{n}^\top \vec{x}=18,\qquad P_-:\vec{n}^\top \vec{x}=14.
\end{align}

Each intersection is a line.  
The direction vector of the line is
\begin{align}
\vec{d} = \vec{m}\times \vec{n} =
\myvec{
\mydet
{m_{23} & n_{23}} \\
 \mydet
 {m_{31} & n_{31}} \\
\mydet
{m_{12} & n_{12}}
}
=
\myvec{1\\0\\-1}.
\end{align}

Pick points:  
For $\Pi\cap P_+$, solve
\begin{align}
\begin{cases}
\vec{m}^\top \vec{x} = 7,\\
\vec{n}^\top \vec{x} = 18
\end{cases}
\quad\Longrightarrow\quad  \vec{p} = \myvec{0\\3\\4}
\end{align}
So
\begin{align}
L_+:\; \vec{r}=\vec{p} + t\vec{d}.
\end{align}

For $\Pi\cap P_-$, solve
\begin{align}
\begin{cases}
\vec{m}^\top \vec{x}= 7,\\
\vec{n}^\top \vec{x}= 14
\end{cases}
\quad\Longrightarrow\quad  \vec{q} = \myvec{0\\7\\0}
\end{align}
So
\begin{align}
L_-:\; \vec{r}=\vec{q} +t\vec{d}.
\end{align}

\textbf{Construct equilateral triangle:}

Take
\begin{align}
\vec{u}=\vec{p}+s\vec{d},\quad \vec{v}=\vec{p}-s\vec{d},\quad \vec{w}=\vec{q}+t\vec{d}.
\end{align}

\begin{align}
{\norm{\vec{u}-\vec{v}}}^2={\norm{\vec{u}-\vec{w}}}^2={\norm{\vec{v}-\vec{w}}}^2.
\end{align}

\begin{align}
\vec{u}-\vec{v}=2s\vec{d} \;\;\Rightarrow\;\; {\norm{\vec{u}-\vec{v}}}^2=4s^2{\norm{\vec{d}}}^2.
\end{align}

\begin{align}
{\norm{\vec{d}}}^2=\vec{d}^\top\vec{d}=1^2+0^2+(-1)^2=2,
{\norm{\vec{u}-\vec{v}}}^2=8s^2.
\end{align}

\begin{align}
\vec{r}=\vec{p}-\vec{q}=\myvec{0\\-4\\4},\quad
\vec{u}-\vec{w}=\vec{r}+(s-t)\vec{d}.
\end{align}

\begin{align}
{\norm{\vec{u}-\vec{w}}}^2={\norm{\vec{r}}}^2+2(s-t)(\vec{r}^\top\vec{d})
+(s-t)^2{\norm{\vec{d}}}^2.
\end{align}

\begin{align}
{\norm{\vec{r}}}^2=32,\quad \vec{r}^\top\vec{d}=-4,\quad {\norm{\vec{d}}}^2=2.
\end{align}

\begin{align}
{\norm{\vec{u}-\vec{w}}}^2=32-8(s-t)+2(s-t)^2.
\end{align}

\begin{align}
\vec{v}-\vec{w}=\vec{r}+(-s-t)\vec{d} \;\;\Rightarrow\;\;
{\norm{\vec{v}-\vec{w}}}^2=32+8(s+t)+2(s+t)^2.
\end{align}

Equating
\begin{align}
32-8(s-t)+2(s-t)^2=8s^2, \quad
32+8(s+t)+2(s+t)^2=8s^2.
\end{align}
Subtracting,
\begin{align}
16s+8st=0 \quad\Rightarrow\quad t=-2,\; s^2=4.
\end{align}

Take $ s=2,t=-2:$
\begin{align}
\vec{u}=\vec{p}+2\vec{d}=\myvec {2\\3\\2},\quad
\vec{v}=\vec{p}-2\vec{d}=\myvec {-2\\3\\6},\quad
\vec{w}=\vec{q}-2\vec{d}=\myvec{-2\\7\\2}.
\end{align}

\textbf{Volume of parallelepiped:}

\textbf{Gram matrix }\(G\) (matrix of pairwise inner products):
\begin{align}
G=\myvec
{{\norm{\vec{u}}}^2 & \vec{u}^\top\vec{v} & \vec{u}^\top\vec{w}\\[6pt]
\vec{v}^\top\vec{u} & {\norm{\vec{v}}}^2 & \vec{v}^\top\vec{w}\\[6pt]
\vec{w}^\top\vec{u} & \vec{w}^\top\vec{v} & {\norm{\vec{w}}}^2}
=
\myvec
{17 & 17 & 21\\[4pt]
17 & 49 & 37\\[4pt]
21 & 37 & 57},
\end{align}


\textbf{Determinant of the Gram matrix (gives the squared volume):}
\begin{align}
\det G =det \myvec
{17 & 17 & 21\\[4pt]
17 & 49 & 37\\[4pt]
21 & 37 & 57} = 12544
\end{align}
Hence the volume itself is
\begin{align}
V=\sqrt{\det G}=\sqrt{12544}=112.
\end{align}
now 
\begin{align}
\frac{80V}{3}= \frac{80*112}{3} = \frac{8960}{3}.
\end{align}
\begin{figure}[h!]
  \centering
  \includegraphics[width=0.9\columnwidth]{figs/fig5.png} 
   \caption*{Fig: Representation of Planes and vectors}
  \label{Fig5}
\end{figure}

\end{document}
